\section{Conclusion}
\label{sec:conclusion}

\emph{Can a table of independently measured primitive costs predict the
behaviour of real composite kernels well enough to drive engineering
decisions?} Where a single resource binds, yes, and tightly: the
per-pipe-max projection sits within 3.5\% on every single-resource
anchor, the smem-bound attention variant within 5\%, the latency-regime
exchange primitive within 10--14\%. Where work couples across issue
pipes, or where a primitive row is a floor rather than a typical cost,
no --- and both failure modes were caught by pre-registered gates rather
than discovered in production. The practical reading for anyone pricing
similar hardware: trust the table's rows, trust compositions in
single-resource regimes, and measure directly anywhere pipes couple.

\subsection{Limitations}

One rig, one toolchain-and-driver snapshot, and rows therefore
confounded with both: a different \texttt{nvcc} emits different SASS for
the same PTX (this paper catalogues seven distinct rewrites), and a
different driver may move the host-domain rows. T4 priors transfer only
core-domain, as recorded per row. Several rows are single-method and
flagged \texttt{UNVERIFIED} pending a second construction; several more
(enumerated in the public coverage manifest) remain open, among them
\texttt{LDSM}, vote, atomic throughputs, and the peer-write visibility
row that blocks the 20~KiB exchange comparison. The chain-method atomics
offset shows that even agreeing contention structure can ride a method
constant. No independent replication exists yet; the repository is
arranged so that one requires a clone and a locked clock.

\subsection{Future work}

The named rows above; a width-aware launch row that would let the
dispatch composition bind; and the same table for the second
architecture this rig's market segment offers cheaply, GA102.
